\documentclass[../algebraIII_notes.tex]{subfiles}
\begin{document}
\section{Aula 06 - 24 de Março, 2026}
\subsection{Motivações}
\begin{itemize}
	\item Polinômios Irredutíveis;
\end{itemize}
\subsection{Polinômios Irredutíveis e Raízes}
\begin{def*}
	Um \textbf{polinômio irredutível sobre F} é um polinômio \(p(x)\in F[x]\) tal que, se
	\[
		p(x)=q_{1}(x)q_{2}(x), \quad q_{1}, q_{2}\in F[x],
	\]
	então ou \(\deg{(q_{1})}=0\), ou \(\deg{(q_{(2)}}=0. \; \square\)
\end{def*}
\begin{lemma*}
	Suponha que F é um corpo e \(p(x)\in F[x]\) é um polinômio irredutível sobre F; então,
	\[
		\frac{F[x]}{\langle p(x) \rangle}
	\]
	é um corpo.
\end{lemma*}
O corpo acima é o das classes de x módulo p, ou seja,
\[
	[x]\in \frac{F[x]}{\langle p(x) \rangle},
\]
com zero sendo
\begin{align*}
	0=[p(x)] & = [a_{0}+a_1x+\dotsc +a_{n}x^{n}]         \\
	         & = [a_{0}]+[a_1x]+ \dotsc + [a_{n}x^{n}]   \\
	         & = [a_{0}]+[a_1][x]+\dotsc +[a_{n}][x^{n}] \\
	         & = a_{0}+a_{1}[x]+\dotsc +a_{n}[x]^{n}.
\end{align*}
\begin{lemma*}
	Se \(p(x)\in F[x]\) é irredutível sobre F, então existe uma extensão de F na qual é possível encontrar uma raiz de \(p(x)\).
\end{lemma*}
\begin{proof*}
	Basta notar que
	\[
		F \subseteq \frac{F[x]}{\langle p(x) \rangle}
	\]
	é uma extensão de F e que \([x]\) é uma raiz do polinômio \(p(x)\in \frac{F[t]}{\langle p(t) \rangle}[x]\)
\end{proof*}

Agora, suponha que K é uma extensão de F e que \(p(x)\in F[x]\) é irredutível, tal que \(\displaystyle \frac{F[x]}{\langle p(x) \rangle}\) é um corpo estendendo F, e seja \(\alpha , \beta \) raízes de \(p(x)\) em K.
\begin{lemma*}
	Existe um único homomorfismo de corpos \( \psi \), com \(\psi (\alpha )=\beta \), tal que:
	\begin{align*}
		 & (1):\; \psi :F(\alpha )\rightarrow F(\beta ) \\
		 & (2):\; \psi (a)=a, \quad \forall a\in F
	\end{align*}
\end{lemma*}
\begin{proof*}
	Consideremos o amigável diagrama trapezoidal
	\begin{center}
		\begin{tikzpicture}[
				observed/.style = {rectangle, thick, text centered, draw, text width = 6em},
				latent/.style = {ellipse, thick, draw, text centered, text width = 6em},
				error/.style ={circle, thick, draw, text centered},
				confounding/.style = {rectangle, thick, text centered, draw, text width = 6em, minimum width = 5.5in},
				outcome/.style = {rectangle, thick, draw, text centered, minimum height = 3.5in, text width = 6em},
			]
			\node(TL) at (-2,1){\(F[x]=F(\alpha )\)};
			\node(BL) at (-2,-2){\(\displaystyle \frac{F[x]}{\langle p(x) \rangle}\)};
			\node(BBL) at (-5,-2){\(F[x]\)};
			\node(TR) at (2,1){\(F(\beta )\)};
			\node(BR) at (2,-2){\(\displaystyle \frac{F[x]}{\langle p(x) \rangle}\)};
			\node(BLB) at (-2,-4){F};
			\node(BRB) at (2,-4){F};
			\node(BBR) at (5,-2){\(F[x]\)};


			\draw[Arrow](BRB)--(BR);
			\draw[<->](BRB)--node[midway, below]{\(= \)}(BLB);
			\draw[Arrow](BLB)--(BL);
			\draw[Arrow, dashed](TL)--node[midway, above]{\(\psi \)}(TR);
			\draw[Arrow](BL)--node[midway, above] {Id}(BR);
			\draw[Arrow](BL)--node[midway, right] {\(\widetilde{\mathrm{ev}_{\alpha }}\)}(TL);
			\draw[Arrow](BR)--node[midway, left] {\(\widetilde{\mathrm{ev}_{\beta }}\)}(TR);
			\draw[Arrow](BBR)--node[midway, above ] {\(\pi \)}(BR);
			\draw[Arrow](BBR)--node[midway, right ] {\(\mathrm{ev}_{\beta } \)}(TR);
			\draw[Arrow](BBL)--node[midway, above ] {\(\pi \)}(BL);
			\draw[Arrow](BBL)--node[midway, left ] {\(\mathrm{ev}_{\alpha } \)}(TL);

		\end{tikzpicture}
	\end{center}
	Por ele, determinamos que \(\psi = \widetilde{\mathrm{ev}_{\beta }}\circ \mathrm{Id}\circ (\widetilde{\mathrm{ev}_{\alpha }})^{-1}\) é um isomorfismo de corpos, onde \(\widetilde{\mathrm{ev}_{\alpha }}([x])=\alpha \), tal que
	\begin{align*}
		\psi (\alpha ) & = \widetilde{\mathrm{ev}_{\beta }} \circ \mathrm{id}\circ (\widetilde{\mathrm{ev}_{\alpha }})^{-1}(\alpha ) \\
		               & = \widetilde{\mathrm{ev}_{\beta }}(\mathrm{id}([x]))                                                        \\
		               & = \widetilde{\mathrm{ev}_{\beta }}([x])                                                                     \\
		               & =\beta
	\end{align*}
	e, para todo a de F,
	\[
		\psi (a)=a.
	\]

	\textbf{\underline{Unicidade}:} considerando
	\[
		F(\alpha )=\{\xi = a_{0}+a_{\alpha } +\dotsc + a_{n-1}\alpha^{n-1}:\; a_{i}\in F\}, \; n = \deg{(p)},
	\]
	então se \(\psi \) é o isomorfismo determinado por
	\[
		\psi (a)=a, \quad \forall a \in F \quad\&\quad \psi (\alpha )=\beta ,
	\]
	segue que, para todo \(\xi \) em \(F(\alpha )\), devemos ter
	\[
		\psi (\xi )=\psi (a_{0}+a_{1}\alpha +\dotsc + a_{n-1}\alpha^{n-1}) = a_{0}+a_{1}\beta +\dotsc + a_{n-1}\beta^{n-1}.
	\]
	Portanto, se \(\psi' :F(\alpha )\rightarrow F(\beta )\) também é um isomorfismo de corpos como acima, então ele satisfaz a mesma identidade para todo \(\xi \) em F, concluindo que
	\[
		\psi (\xi )=\psi '(\xi ),\quad \forall \xi \in F \Rightarrow \psi \equiv \psi '. \text{ \qedsymbol}
	\]
\end{proof*}
\begin{crl*}
	Sejam K uma extensão de corpo de F, \(p(x)\in F[x]\) irredutível e \(\alpha \in K\) uma raiz fixada de \(p(x)\). Então, existe uma correspondência bijetora entre os conjuntos
	\[
		\biggl\{\substack{\text{Raízes de }p(x) \\ \text{em K}}\biggr\}
		\longleftrightarrow \biggl\{\substack{\eta :F(\alpha )\rightarrow K \text{ homomorfismo de} \\ \text{corpos tal que }\eta (a)=a,\; \forall a\in F }\biggr\}
	\]
\end{crl*}
\begin{exr}
	Prove o corolário acima.
\end{exr}
\begin{lemma*}
	Consideremos duas extensões de dois corpos, \(F\subseteq K,\; F'\subseteq K'\), um polinômio \(p(x)\in F[x]\) irredutível, um isomorfismo de corpos \(\varphi :F\rightarrow F'\) e o polinômio \(p'(x)\coloneqq (\varphi_{*}p)(x)\in F'[x]\).
	Suponha que \(\alpha \in K\) é raiz de \(p(x)\) e que \(\alpha '\in K'\) é raiz de m  \(p'(x)\); então, existe um isomorfismo de corpos único \(\psi :F(\alpha )\rightarrow F'(\alpha ')\) tal que:
	\begin{align*}
		 & (1):\; \psi (\alpha )=\alpha '                     \\
		 & (2):\; \psi (a )=\varphi(a), \quad \forall a\in F.
	\end{align*}
\end{lemma*}
\begin{proof*}
	Escrevamos o seguinte diagrama:
	\begin{center}
		\begin{tikzpicture}[
				observed/.style = {rectangle, thick, text centered, draw, text width = 6em},
				latent/.style = {ellipse, thick, draw, text centered, text width = 6em},
				error/.style ={circle, thick, draw, text centered},
				confounding/.style = {rectangle, thick, text centered, draw, text width = 6em, minimum width = 5.5in},
				outcome/.style = {rectangle, thick, draw, text centered, minimum height = 3.5in, text width = 6em},
			]
			\node(TL) at (-1,1){\(\displaystyle \frac{F[x]}{\langle p(x) \rangle}\)};
			\node(BL) at (-1,-1){F};
			\node(TR) at (1,1){\(\displaystyle \frac{F'[x]}{\langle p'(x) \rangle}\)};
			\node(BR) at (1,-1){F};

			\draw[Arrow](TL)--node[midway, above] {\(\tilde{\phi}\)}(TR);
			\draw[Arrow](BL)--node[midway, below] {\(\phi\)}(BR);
			\draw[Arrow](TL)--node[midway, left] {}(BL);
			\draw[Arrow](TR)--node[midway, right] {}(BR);

		\end{tikzpicture}
	\end{center}
	A partir dele, definimos
	\[
		\left\{\begin{array}{ll}
			\tilde{\varphi }(a)=\varphi (a),\; \forall a\in F \\
			\tilde{\varphi }([x]_{p(x)})=[x]_{p'(x)}
		\end{array}\right.,
	\]
	tal que
	\[
		[x]_{p(x)} = x + \langle p(x) \rangle,
	\]
	permitindo fecharmos o diagrama por cima com
	\begin{center}
		\begin{tikzpicture}[
				observed/.style = {rectangle, thick, text centered, draw, text width = 6em},
				latent/.style = {ellipse, thick, draw, text centered, text width = 6em},
				error/.style ={circle, thick, draw, text centered},
				confounding/.style = {rectangle, thick, text centered, draw, text width = 6em, minimum width = 5.5in},
				outcome/.style = {rectangle, thick, draw, text centered, minimum height = 3.5in, text width = 6em},
			]
			\node(TTL) at (-1,3){\(F(\alpha )\)};
			\node(TTR) at (1,3){\(F(\beta )\)};
			\node(TL) at (-1,1){\(\displaystyle \frac{F[x]}{\langle p(x) \rangle}\)};
			\node(BL) at (-1,-1){F};
			\node(TR) at (1,1){\(\displaystyle \frac{F'[x]}{\langle p'(x) \rangle}\)};
			\node(BR) at (1,-1){F};

			\draw[Arrow, dashed](TTL)--node[midway, above] {\(\psi \)}(TTR);
			\draw[Arrow](TL)--node[midway, above] {\(\tilde{\phi}\)}(TR);
			\draw[Arrow](BL)--node[midway, below] {\(\phi\)}(BR);
			\draw[Arrow](BL)--node[midway, left] {}(TL);
			\draw[Arrow](BR)--node[midway, right] {}(TR);
			\draw[Arrow](TL)--node[midway, left] {\(\widetilde{\mathrm{ev}_{\alpha }}\)}(TTL);
			\draw[Arrow](TR)--node[midway, right] {\(\widetilde{\mathrm{ev}_{\alpha '}}\)}(TTR);
		\end{tikzpicture}
	\end{center}
	Seguindo os caminhos do diagrama, determinamos as relações que provam o teorema:
	\begin{align*}
		 & \tikz[baseline=-0.5ex]\draw[black, fill=black, radius=1.5pt](0, 0)circle;\; \psi = \widetilde{\mathrm{ev}_{\alpha '}}\circ \tilde{\phi}\circ (\widetilde{\mathrm{ev}_{\alpha }})^{-1}                                                \\
		 & \tikz[baseline=-0.5ex]\draw[black, fill=black, radius=1.5pt](0, 0)circle;\; \psi (\alpha )= \widetilde{\mathrm{ev}_{\alpha '}}\circ \tilde{\phi}\circ (\widetilde{\mathrm{ev}_{\alpha }})^{-1}(\alpha )=\alpha '. \text{ \qedsymbol}
	\end{align*}
\end{proof*}
\begin{tcolorbox}[
		skin=enhanced,
		title=Observação,
		fonttitle=\bfseries,
		colframe=black,
		colbacktitle=cyan!75!white,
		colback=cyan!15,
		colbacklower=black,
		coltitle=black,
		drop fuzzy shadow,
		%drop large lifted shadow
	]
	Em particular, o lema acima dá a relação que \(p(x)\) é irredutível se, e somente se, \(p'(x)=(\varphi_{*}p)(x)\) for irredutível.
\end{tcolorbox}
Usando estes lemas e resultados, nosso objetivo é, dado \(f(x)\) definir o que chamamos de grupo de Galois de \(f(x)\) a partir da extensão \(K_{f}\) como uma extensão de corpo de polinômios e os isomorfismos de corpos \(\varphi :K_{f}\rightarrow K_{f}\) tais que \(\varphi (a)=a\) para todo a em F, que resultará no corpo de decomposição de f.
O grupo de Galois de \(f(x)\) será exatamente o grupo das funções \(\varphi \) mencionadas!

\subsection{Corpo de Decomposição de um Polinômio}
Vamos considerar, adiante, a notação fixa para \(f(x)\in F[x]\), onde F é um corpo
\begin{def*}
	Seja K um corpo com \(F\subseteq K\); então, K é chamado \textbf{corpo de decomposição de \(\mathbf{f(x)}\) sobre F} se:
	\begin{itemize}
		\item[1)] o polinômio \(f(x)\) pode ser escrito em termos de polinômios minimais:
		      \[
			      f(x)=a \prod\limits_{j=1}^{\deg{(f)}}(x-\alpha_{j}),
		      \]
		      onde \(a\in F\) e \(\alpha_{j}\in K\) para \(j=1,\dotsc , \deg{(f)}\);
		\item[2)] K é \textit{minimal} em relação a esta propriedade, ou seja, se \(K'\) é um outro corpo que contém F, satisfaz
		      \[
			      f(x)=b \prod\limits_{j=1}^{\deg{(f)}}(x-\beta_{j}),\quad b\in F, \; \beta_{j}\in K'
		      \]
		      e \(K'\subseteq K\) propriamente, então \(K=K',\;\alpha_{j}=\beta_{j}\) e \(a=b\).
	\end{itemize}
\end{def*}
\begin{example}
	Vamos considerar o caso em que \(F=\mathbb{Q}\) e \(f(x)=x^{2}-2\in \mathbb{Q}[x]\); então, \(\mathbb{Q}(\sqrt[]{2})\) é um corpo de decomposição de \(f(x)|_{\mathbb{Q}}\), assim como \(\displaystyle \frac{\mathbb{Q}[x]}{\langle x^{2}-2 \rangle}\).

	Note que eles não são iguais, mas são isomorfos pelo isomorfismo
	\[
		\widetilde{\mathrm{ev}_{\alpha }}:\frac{\mathbb{Q}[x]}{\langle x^{2}-2 \rangle}\rightarrow \mathbb{Q}(\sqrt[]{2}).
	\]
\end{example}
\begin{example}
	Para \(F=\mathbb{Q}\) e \(f(x)=(x^{2}-2)(x^{2}-3)\), os polinômios \(h(x)=x^{2}-2\) e \(g(x)=x^{2}-3\) são ambos irredutíveis sobre \(\mathbb{Q}\), tal que o candidato a corpo de decomposição é
	\[
		K=\mathbb{Q}(\sqrt[]{2}, \sqrt[]{3}),
	\]
	pois é nele que encontramos todas as raízes de \(h, \; f\) e g, e é o menor corpo que contém essa propriedade. Observe também que
	\[
		\mathbb{Q}(\sqrt[]{2}, \sqrt[]{3})=\mathbb{Q}(\sqrt[]{2}+\sqrt[]{3}),
	\]
	tal que se demonstrarmos essa observação, ganhamos o grau da extensão como 4, conforme mostramos anteriormente.

	Um lado dessa demonstração é fácil, bastando notar que o elemento \(\sqrt[]{2}+\sqrt[]{3}\) é um elemento de \(\mathbb{Q}(\sqrt[]{2}, \sqrt[]{3})\); por outro lado, temos que provar que \(\sqrt[]{2},; \sqrt[]{3}\) pertencem a \(\mathbb{Q}(\sqrt[]{2}+\sqrt[]{3})\). Com efeito, note que
	\[
		(\sqrt[]{2}+\sqrt[]{3})^{2}=5+2\sqrt[]{6},
	\]
	mostrando que \(\sqrt[]{6}\) pertence a \(\mathbb{Q}(\sqrt[]{2}+\sqrt[]{3})\); assim,
	\[
		\sqrt[]{6}(\sqrt[]{2}+\sqrt[]{3}) = 2\sqrt[]{3}+3\sqrt[]{2}\in \mathbb{Q}(\sqrt[]{2}+\sqrt[]{3}).
	\]
	Agora, podemos manipular algebricamente com o termo \(-2(\sqrt[]{2}+\sqrt[]{3})+2\sqrt[]{3}+3\sqrt[]{2}\) para obter a expressão
	\[
		2\sqrt[]{3}+3\sqrt[]{2}-2(\sqrt[]{2}+\sqrt[]{3})+2\sqrt[]{3}+3\sqrt[]{2} \Rightarrow \sqrt[]{2}\in \mathbb{Q}(\sqrt[]{2}+\sqrt[]{3}).
	\]
	Analogamente, prova-se que \(\sqrt[]{3}\in \mathbb{Q}(\sqrt[]{2}+\sqrt[]{3})\).

	Portanto, concluímos que
	\[
		[\mathbb{Q}(\sqrt[]{2}+\sqrt[]{3}):\mathbb{Q}]=4,
	\]
	que é o mesmo grau de \(f(x)\).
\end{example}
\begin{example}
	Dessa vez, vamos tomar um polinômio de grau 3 analisando o caso \(F=\mathbb{Q}\) e \(f(x)=x^{3}-2\in \mathbb{Q}[x]\), que é irredutível sobre \(\mathbb{Q}\).
	Podemos considerar o quociente
	\[
		\frac{\mathbb{Q}[x]}{\langle x^{3}-2 \rangle}\cong \mathbb{Q}(\sqrt[3]{2}),
	\]
	que fornece uma extensão de grau 3. O problema é que este corpo não é o corpo de decomposição de \(f(x)\), pois nem todas as raízes estão ali dentro -- uma raiz é real, que é \(\sqrt[3]{2}\), e duas raízes complexas, que consequentemente não estão no corpo.
	Consequentemente, este corpo de fato não é um corpo de decomposição de \(f(x)\).
\end{example}
\begin{exr}
	Mostre que o corpo de decomposição da \(f(x)\) do exemplo acima é \(\mathbb{Q}(\sqrt[3]{2}, \sqrt[]{-3})\) e que
	\[
		[\mathbb{Q}(\sqrt[3]{2}, \sqrt[]{-3}): \mathbb{Q}]=6.
	\]
\end{exr}

\end{document}
